% Book-owned openings and handoffs. The seven chapter sources remain
% independently buildable; these rails exist only in the collected volume.

\newcommand{\pdbookrail}[2]{%
  \par\vspace{0.5cm}%
  \noindent{\color{hhcobalt}\rule{2.5em}{0.8pt}}\par
  \vspace{0.12cm}%
  \noindent{\small\scshape\bfseries\color{hhcobalt}#1}\enspace
  {\small #2}\par
  \vspace{0.45cm}%
}

\newcommand{\pdchapteropeningI}{%
  \pdbookrail{Question I: what must remain visible?}{%
    A swarm can produce more events than its operator can read. The resulting
    scarcity is not cured by a smoother summary: some summaries make the
    decisive artifact unreachable. This chapter establishes the cost of a
    zero-miss digest, proves why a successor and an operator require different
    rankings, and derives the attention rule from explicit error costs. Its
    strongest claims are governed by R1--R4, R14, and R16.}%
}

\newcommand{\pdchapterhandoffI}{%
  \pdbookrail{What Chapter II receives}{%
    The operator now has a disciplined read surface: claims remain linked to
    artifacts, attention is priced, and uncertainty is visible. None of that
    can stop a write. The next chapter moves from what can be seen to the single
    local boundary where a proposed action can become a durable effect.}%
}

\newcommand{\pdchapteropeningII}{%
  \pdbookrail{Question II: where can a rule be made real?}{%
    Six agents may agree about a file and still overwrite one another. A log can
    report the collision after the fact; only a mediated writer can refuse the
    second effect before it lands. This chapter identifies that local writer,
    separates preventable rules from detect-only claims, and states the work-unit
    invariants that every later institution assumes. R5 and R8 set the boundary.}%
}

\newcommand{\pdchapterhandoffII}{%
  \pdbookrail{What Chapter III receives}{%
    Local effects now have one serial history, one fencing epoch, and one place
    to attach receipts. The writer does not solve continuity: a process can die,
    restart under a new model, or shed its local memory. The next chapter asks
    what must persist so yesterday's act can still be attributed tomorrow.}%
}

\newcommand{\pdchapteropeningIII}{%
  \pdbookrail{Question III: what remains long enough to owe anything?}{%
    Accountability cannot attach to a disposable process name. It needs a
    durable identity, a named principal, an outcome history, and rules for
    migration that do not erase sanctions or mint reputation. This chapter
    builds that continuity without pretending software has a human soul. R12,
    R13, and B6 govern its economic and protocol claims.}%
}

\newcommand{\pdchapterhandoffIII}{%
  \pdbookrail{What Chapter IV receives}{%
    A successor can now inherit a bounded history without inheriting imaginary
    credit or escaping old obligations. Continuity by itself creates no market:
    it neither prices risk nor says whose funds may move. The next chapter turns
    attributed work into exchange by separating value, collateral, evidence,
    and settlement.}%
}

\newcommand{\pdchapteropeningIV}{%
  \pdbookrail{Question IV: what makes cooperation rational without making loss unbounded?}{%
    A payment rail moves money; a market must also identify the parties, state
    what performance means, reserve the right funds, and decide what happens
    when evidence disagrees. This chapter constructs that market from four
    separate monetary buckets and explicit oracles. R7 and R13--R15 delimit the
    audit, substitution, escalation, and specialization claims.}%
}

\newcommand{\pdchapterhandoffIV}{%
  \pdbookrail{What Chapter V receives}{%
    Work can now be proposed, bonded, graded, and settled against a named
    principal. The missing question is prior to price: why may this actor invoke
    this effect at all, and what stops delegated authority from growing at the
    next hop? The next chapter turns authority into a key-bound, attenuating,
    revocable object.}%
}

\newcommand{\pdchapteropeningV}{%
  \pdbookrail{Question V: how may authority travel without silently growing?}{%
    Identity answers who signed; it does not answer what the signer may do.
    Anchor therefore checks every delegation edge, binds the grant to a key and
    audience, narrows scope and lifetime, and preserves revocation ancestry.
    The formal claim concerns the checked protocol and Rust paths, while R5
    explains which surrounding policy promises can become admission gates.}%
}

\newcommand{\pdchapterhandoffV}{%
  \pdbookrail{What Chapter VI receives}{%
    A request now arrives with authenticated, narrowing authority. That still
    does not make its result true, its evidence sufficient, or its remedy fair.
    The next chapter composes grants with witnessed history, conflict checks,
    audit, appeals, and one bounded settlement.}%
}

\newcommand{\pdchapteropeningVI}{%
  \pdbookrail{Question VI: who may decide, on what evidence, with what consequence?}{%
    Attribution without consequence is only a diary; consequence without
    admissible evidence is arbitrary. This chapter turns the local writer into
    an institution by joining capability bounds, witnessed records, audits,
    appeals, and conserved settlement. R7--R11, R15, R17, and the R8 machine
    state exactly what has been proved, checked, or left conditional.}%
}

\newcommand{\pdchapterhandoffVI}{%
  \pdbookrail{What Chapter VII receives}{%
    One machine can now authorize, observe, adjudicate, and settle its own work.
    A second machine introduces a different operator and a different source of
    truth. The final chapter asks which capabilities and records may cross that
    boundary without pretending that either local authority has become global.}%
}

\newcommand{\pdchapteropeningVII}{%
  \pdbookrail{Question VII: what can cross a trust boundary without moving the authority itself?}{%
    The current Relay federates harbor-scoped events over outbound HTTPS and
    SSE; it does not replicate daemon state or create consensus between
    operators. Local admission remains local. The research result is narrower
    and sharper still: R6 and CR-1/2/3 show exactly when relayed reports around a
    cycle can expose an inconsistency that a missing pairwise check cannot.}%
}

\newcommand{\pdchapterhandoffVII}{%
  \pdbookrail{The conclusion of the book}{%
    The seven answers now form one discipline. Preserve enough evidence for a
    human to challenge the system; mediate every effect that must be prevented;
    attach actions to a durable principal and work unit; let delegated authority
    narrow; make settlement conserve value; and cross organizational boundaries
    by exchanging verifiable records rather than by exporting sovereignty.
    Better models help inside this discipline. They do not replace it.}%
}
